\subsubsection*{Solution}

Take the common polarization to be \(\hat{\mathbf y}\), and choose the
particle labels so that

\[
\mathbf r_1-\mathbf r_2=(d_0,0,z_1-z_2).
\]

The positive quantities \(E_1\) and \(E_2\) denote the focal-plane field
amplitudes, and \(\alpha_1,\alpha_2\in\mathbb C\) are the physical scalar
polarizabilities at the laser frequency.  We use the phasor convention

\[
\mathbf E(\mathbf r,t)=\operatorname{Re}\!\left[\mathbf E(\mathbf r)e^{-i\omega t}\right].
\]

The displacements are measured from the actual static equilibrium.  In the
equations of motion below, the term \(-m\Omega_j^2z_j\) includes every local
linear force proportional only to \(z_j\), including the axial
intensity-curvature correction.  The coefficients \(k_1\) and \(k_2\)
describe the remaining relative-coordinate optical-binding force.

\paragraph*{Incident fields near the focal planes}

On the propagation axis, the field of tweezer $j$ is

\[
\mathbf E_j(z_j) = \hat{\mathbf y}\, \frac{E_j}{\sqrt{1+(z_j/z_R)^2}} \exp\!\left[ i\left(kz_j-\arctan\frac{z_j}{z_R}+\phi_j\right) \right].
\]

Because $|z_j|\ll z_R$,

\[
\frac{1}{\sqrt{1+(z_j/z_R)^2}} =1+O\!\left(\frac{z_j^2}{z_R^2}\right),
\]

and

\[
\arctan\frac{z_j}{z_R} = \frac{z_j}{z_R} +O\!\left(\frac{z_j^3}{z_R^3}\right).
\]

Thus, to linear order,

\[
\mathbf E_j(z_j) \simeq \hat{\mathbf y}\,E_j e^{i(qz_j+\phi_j)}, \qquad q\equiv k-\frac{1}{z_R}.
\]

Consequently, the local relative phase is

\[
\Delta\phi = \phi_1-\phi_2+q(z_1-z_2).
\]

The amplitude correction begins at quadratic order in \(z_j\).  Its gradient
therefore produces only a local term proportional to \(z_j\), already included
in \(\Omega_j\).  The relative-coordinate coupling at the retained order is
generated by the phase \(q(z_1-z_2)\).

\paragraph*{Dipole force to bilinear order}

Let \(\mathbf E_j\) denote the incident tweezer field at particle \(j\), and
let \(\mathbf E_{\mathrm{exc},j}\) be the field that excites its dipole.  With
\(\ell\ne j\),

\[
\mathbf E_{\mathrm{exc},j}
=\mathbf E_j+\frac{1}{\varepsilon_0}\mathbf G_{j\ell}\mathbf p_\ell,
\qquad
\mathbf p_j=\alpha_j\mathbf E_{\mathrm{exc},j},
\]

where

\[
\mathbf G_{j\ell}=\mathbf G(\mathbf r_j-\mathbf r_\ell).
\]

In the single-scattering approximation,

\[
\begin{aligned}
\mathbf E_{\mathrm{exc},j}
&=\mathbf E_j+
\frac{\alpha_\ell}{\varepsilon_0}\mathbf G_{j\ell}\mathbf E_\ell,\\
\mathbf p_j
&=\alpha_j\mathbf E_j+
\frac{\alpha_j\alpha_\ell}{\varepsilon_0}
\mathbf G_{j\ell}\mathbf E_\ell,
\end{aligned}
\]

up to higher multiple-scattering orders.  The time-averaged electric-dipole
force is

\[
(F_j)_a=\frac12\operatorname{Re}\!\left[
\mathbf p_j^*\cdot\partial_{j,a}\mathbf E_{\mathrm{exc},j}
\right],
\]

where \(a\in\{x,y,z\}\) and
\(\partial_{j,a}\equiv\partial/\partial(\mathbf r_j)_a\), with the other
particle coordinate held fixed.  Substitution of the preceding
single-scattering fields and retention of the terms bilinear in the two
polarizabilities gives

\[
\begin{aligned}
(F_j^{\mathrm{bind}})_a
=\frac{1}{2\varepsilon_0}\operatorname{Re}\!\Big[{}
&\alpha_j^*\alpha_\ell\,
\mathbf E_j^* \cdot
\partial_{j,a}\!\left(\mathbf G_{j\ell}\mathbf E_\ell\right)\\
&+\alpha_j^*\alpha_\ell^*\,
(\mathbf G_{j\ell}\mathbf E_\ell)^* \cdot
\partial_{j,a}\mathbf E_j
\Big],
\qquad \ell\neq j,
\end{aligned}
\]

In the first term, \(\partial_{j,a}\) acts only on
\(\mathbf G_{j\ell}\), because \(\mathbf E_\ell\) is evaluated at the other
particle.

\paragraph*{Leading far-field reduction}

For \(kd\gg1\), the free-space Green tensor has the leading form

\[
\mathbf G(\mathbf d)
=\frac{k^2e^{ikd}}{4\pi d}
\left(\mathbf I-\hat{\mathbf d}\hat{\mathbf d}\right)+O(d^{-2}).
\]

The separation vector remains in the \(xz\)-plane and is perpendicular to
the common \(y\)-polarization.  Hence the relevant scalar Green component is

\[
g(d)\equiv G_{yy}(d)
=\frac{k^2e^{ikd}}{4\pi d}+O(d^{-2}),
\qquad
d=\sqrt{d_0^2+(z_1-z_2)^2}.
\]

\Needspace{7\baselineskip}
With \(\delta z\equiv z_1-z_2\),

\[
d=\sqrt{d_0^2+\delta z^2} =d_0+O(\delta z^2).
\]

Differentiation gives

\[
\partial_{1,z}d=\frac{\delta z}{d_0}+O(\delta z^3),
\qquad
\partial_{2,z}d=-\frac{\delta z}{d_0}+O(\delta z^3).
\]

Therefore the first term of the binding-force formula produces a linear
coefficient with far-field scaling \(d_0^{-2}\).  The second term contains a
phase-gradient coefficient proportional to \(g(d_0)q^2\), whose leading
far-field scaling is \(d_0^{-1}\).  The calculation retains only the
\(d_0^{-1}\) part of the linearized binding stiffness and omits all
\(d_0^{-2}\) and higher-order far-field corrections.  It therefore uses

\[
g(d_0)=\frac{k^2e^{ikd_0}}{4\pi d_0}
\]

in the phase-gradient force term.

\paragraph*{Linearization and equations of motion}

Let

\[
\Delta\phi_0\equiv\phi_1-\phi_2, \qquad \Delta\phi=\Delta\phi_0+q\delta z.
\]

Using \(\partial_{j,z}\mathbf E_j=iq\mathbf E_j\) to the order relevant for
the relative-coordinate coupling, the leading phase-gradient contribution to
the first binding force is

\[
\begin{aligned}
F_{1z}^{\mathrm{bind}}
={}&\frac{qE_1E_2}{2\varepsilon_0}
\operatorname{Im}\!\left[
\alpha_1\alpha_2g(d_0)e^{-i(\Delta\phi_0+q\delta z)}
\right],
\end{aligned}
\]

and the corresponding contribution to the second force is

\[
\begin{aligned}
F_{2z}^{\mathrm{bind}}
={}&\frac{qE_1E_2}{2\varepsilon_0}
\operatorname{Im}\!\left[
\alpha_1\alpha_2g(d_0)e^{i(\Delta\phi_0+q\delta z)}
\right].
\end{aligned}
\]

Expanding the phases to first order in \(\delta z\) gives

\[
\begin{aligned}
F_{1z}^{\mathrm{bind}}
={}&F_{1z}^{(0)}
-
\frac{E_1E_2k^2q^2}{8\pi\varepsilon_0d_0}
\operatorname{Re}\!\left[
\alpha_1\alpha_2e^{i(kd_0-\Delta\phi_0)}\right]
(z_1-z_2),
\end{aligned}
\]

\[
\begin{aligned}
F_{2z}^{\mathrm{bind}}
={}&F_{2z}^{(0)}
+
\frac{E_1E_2k^2q^2}{8\pi\varepsilon_0d_0}
\operatorname{Re}\!\left[
\alpha_1\alpha_2e^{i(kd_0+\Delta\phi_0)}\right]
(z_1-z_2).
\end{aligned}
\]

The constants \(F_{jz}^{(0)}\) are canceled by the static force balance.
Writing the remaining dynamics as

\[
\begin{aligned}
m\ddot z_1&=-m\Omega_1^2z_1-(k_1+k_2)(z_1-z_2),\\
m\ddot z_2&=-m\Omega_2^2z_2+(k_1-k_2)(z_1-z_2),
\end{aligned}
\]

comparison of the coefficients of \(z_1-z_2\) gives

\[
\begin{aligned}
k_1+k_2={}&
\frac{E_1E_2k^2q^2}{8\pi\varepsilon_0d_0}
\operatorname{Re}\!\left[
\alpha_1\alpha_2e^{i(kd_0-\Delta\phi_0)}\right],
\end{aligned}
\]

and

\[
\begin{aligned}
k_1-k_2={}&
\frac{E_1E_2k^2q^2}{8\pi\varepsilon_0d_0}
\operatorname{Re}\!\left[
\alpha_1\alpha_2e^{i(kd_0+\Delta\phi_0)}\right].
\end{aligned}
\]

Adding and subtracting these two equations, and then substituting
$q=k-1/z_R$, yields

\[
\boxed{\begin{aligned}
k_1={}&\frac{E_1E_2k^2}
{8\pi\varepsilon_0d_0}
\left(k-\frac{1}{z_R}\right)^2
\operatorname{Re}\!\left[
\alpha_1\alpha_2e^{ikd_0}\right]\cos\Delta\phi_0,
\end{aligned}}
\]

\[
\boxed{\begin{aligned}
k_2={}&\frac{E_1E_2k^2}
{8\pi\varepsilon_0d_0}
\left(k-\frac{1}{z_R}\right)^2
\operatorname{Im}\!\left[
\alpha_1\alpha_2e^{ikd_0}\right]\sin\Delta\phi_0.
\end{aligned}}
\]

Both coefficients have units of \(\mathrm{N\,m^{-1}}\).  They are the
\(d_0^{-1}\) terms in the far-field expansion; the omitted directional
stiffness begins at far-field order \(d_0^{-2}\).
